\chapter{Obligations convertibles et lien crédit-action}\label{ch:dv:convertibles-and-the-credit-equity-link}

Une société dont l'action cote 30 émet une obligation à cinq ans que le porteur peut échanger contre 2.5 actions à tout moment. Avec l'action
à 80, l'obligation vaut 202, à peine plus que les 200 que rapporteraient ses actions~: elle se comporte comme l'action. Avec l'action à
10, elle vaut 69, à peine plus que les 68 d'une obligation ordinaire de la même société, et cette obligation ordinaire a elle-même baissé, parce qu'une
société dont l'action s'est effondrée a davantage de risques de faire défaut. Elle se négocie sur son crédit. Entre les deux, l'\omterm{def:dv:convertibles-and-the-credit-equity-link:cb}{obligation convertible} relève à la fois des deux,
et un desk de convertibles évolue dans cette zone. Il détient l'obligation, vend l'action en couverture, achète une protection de crédit et capte la
convexité. Ce chapitre décrit l'instrument, son \omterm{def:dv:convertibles-and-the-credit-equity-link:floor}{plancher obligataire}, sa \omterm{def:dv:convertibles-and-the-credit-equity-link:value}{valeur de conversion} et sa zone convexe centrale. Il construit ensuite un
\omterm{def:dv:convertibles-and-the-credit-equity-link:e2c}{modèle action-crédit} sur une grille aux différences finies, déduit ce que le desk couvre, et se termine sur ce qui s'est passé lorsque le financement de l'arbitrage
a disparu.

\section{L'instrument}

\begin{definition}[Obligation convertible]\label{def:dv:convertibles-and-the-credit-equity-link:cb}
Une \emph{obligation convertible}\index{obligation convertible} (convertible bond) est une obligation que son porteur peut échanger contre un nombre fixe d'actions de l'émetteur,
à tout moment ou pendant des périodes déterminées, au lieu de recevoir les coupons et le principal~; elle confère en général aussi à l'émetteur le droit de la
racheter par anticipation, et parfois au porteur le droit d'en exiger le remboursement.
\end{definition}

\begin{definition}[Ratio de conversion]\label{def:dv:convertibles-and-the-credit-equity-link:ratio}
Le \emph{ratio de conversion}\index{ratio de conversion} (conversion ratio) est le nombre d'actions reçues par obligation lors de la conversion.
\end{definition}

\begin{definition}[Prix de conversion]\label{def:dv:convertibles-and-the-credit-equity-link:price}
Le \emph{prix de conversion}\index{prix de conversion} (conversion price) est la valeur nominale divisée par le \omterm{def:dv:convertibles-and-the-credit-equity-link:ratio}{ratio de conversion}~: le cours de l'action auquel
la conversion vaut la valeur nominale.
\end{definition}

\begin{definition}[Valeur de conversion]\label{def:dv:convertibles-and-the-credit-equity-link:value}
La \emph{valeur de conversion}\index{valeur de conversion} (conversion value, parité) est le \omterm{def:dv:convertibles-and-the-credit-equity-link:ratio}{ratio de conversion} multiplié par le cours actuel de l'action, ce que vaudrait l'obligation
si elle était convertie maintenant.
\end{definition}

\begin{definition}[Prime de conversion]\label{def:dv:convertibles-and-the-credit-equity-link:premium}
La \emph{prime de conversion}\index{prime de conversion} (conversion premium) est le rapport du prix de l'\omterm{def:dv:convertibles-and-the-credit-equity-link:cb}{obligation convertible} à sa \omterm{def:dv:convertibles-and-the-credit-equity-link:value}{valeur de conversion}, moins un.
\end{definition}

L'obligation du chapitre a une valeur nominale de 100, cinq ans à courir, un coupon annuel de 2~\%, et un \omterm{def:dv:convertibles-and-the-credit-equity-link:ratio}{ratio de conversion} de 2.5, soit un \omterm{def:dv:convertibles-and-the-credit-equity-link:price}{prix
de conversion} de 40 pour une action à 30. L'émetteur peut la racheter à 100 à partir de la fin de la deuxième année, mais seulement si l'action vaut au moins 52,
130~\% du \omterm{def:dv:convertibles-and-the-credit-equity-link:price}{prix de conversion}. Les taux sont à 3~\%, le rendement du dividende à 1~\% et la volatilité de l'action à 30~\%. Le spread de crédit de l'émetteur est de 300
points de base, une intensité de défaut de 5~\% par an avec un recouvrement de 40~\% de la valeur nominale.

\begin{definition}[Soft call]\label{def:dv:convertibles-and-the-credit-equity-link:softcall}
Un \emph{soft call}\index{soft call} est le droit de l'émetteur de rembourser par anticipation une \omterm{def:dv:convertibles-and-the-credit-equity-link:cb}{obligation convertible}, exerçable seulement si le cours de l'action
a dépassé un seuil déclencheur (par exemple 130~\% du \omterm{def:dv:convertibles-and-the-credit-equity-link:price}{prix de conversion})~: un call contraint les porteurs
à convertir lorsque la conversion vaut plus que le prix de rachat.
\end{definition}

Le call de l'émetteur est un coût pour le porteur. Le rachat force la conversion lorsque la \omterm{def:dv:convertibles-and-the-credit-equity-link:value}{valeur de conversion} dépasse le prix de rachat. Cela plafonne la
valeur de l'obligation près du seuil et supprime la valeur temps que le porteur conserverait autrement. Dans le modèle ci-dessous, le \omterm{def:dv:convertibles-and-the-credit-equity-link:softcall}{soft call} vaut
3.9 pour un cours de l'action de 40 et 4.7 pour 52~: 115.3 et 138.6 avec le call, contre 119.2 et 143.3 sans lui.

\section{Plancher obligataire, valeur de conversion et zone convexe}

\begin{definition}[Plancher obligataire]\label{def:dv:convertibles-and-the-credit-equity-link:floor}
Le \emph{plancher obligataire}\index{plancher obligataire} (bond floor) d'une convertible est la valeur de la même obligation sans le droit de conversion~: ses coupons et
le principal, actualisés au risque de crédit de l'émetteur (livre~2, chapitre~21).
\end{definition}

Une \omterm{def:dv:convertibles-and-the-credit-equity-link:cb}{obligation convertible} vaut au moins son \omterm{def:dv:convertibles-and-the-credit-equity-link:floor}{plancher obligataire}, car le porteur peut simplement garder l'obligation, et au moins sa \omterm{def:dv:convertibles-and-the-credit-equity-link:value}{valeur de conversion},
car elle peut se convertir. Elle vaut plus que chacune des deux, de la valeur de l'option de choisir plus tard. Loin du \omterm{def:dv:convertibles-and-the-credit-equity-link:price}{prix de conversion}, l'une
des deux bornes domine et la convertible se comporte comme cette borne~: une obligation à gauche, les actions à droite. Près du \omterm{def:dv:convertibles-and-the-credit-equity-link:price}{prix de conversion}, les deux
comptent. Là, la convertible est convexe par rapport à l'action, et c'est cette convexité que recherche l'arbitrage
(la \cref{fig:dv:convertibles-and-the-credit-equity-link:profile}).

\begin{example}[La convertible selon les cours de l'action]\label{ex:dv:convertibles-and-the-credit-equity-link:profile}
Sous une intensité de défaut constante de 5~\%, la convertible vaut 89.1, 100.2, 116.0, 156.7 et 203.2 pour des cours de l'action de 20, 30, 40, 60 et 80. Son \omterm{def:dv:convertibles-and-the-credit-equity-link:floor}{plancher
obligataire} est de 83.2 quel que soit le cours. Les \omterm{def:dv:convertibles-and-the-credit-equity-link:value}{valeurs de conversion} sont de 50, 75, 100, 150 et 200, si bien que la \omterm{def:dv:convertibles-and-the-credit-equity-link:premium}{prime de conversion} passe de 78~\% au cours de 20 à 34~\% au cours de 30,
16~\% au cours de 40, 4.5~\% au cours de 60 et 1.6~\% au cours de 80. Le delta passe de 0.84 actions par obligation au cours de 20 à 2.39 au cours de 80, se rapprochant du \omterm{def:dv:convertibles-and-the-credit-equity-link:ratio}{ratio de conversion} de 2.5.
\end{example}

\begin{omfigure}
\begin{tikzpicture}
\begin{axis}[width=10.5cm,height=5.4cm,
  xlabel={cours de l'action},ylabel={valeur pour 100 de nominal},
  tick label style={font=\small},label style={font=\small},
  xmin=10,xmax=90,ymin=0,ymax=230,grid=major,grid style={gray!25},
  legend columns=2,legend cell align=left,legend style={font=\footnotesize,at={(0.5,-0.28)},anchor=north,draw=none,fill=none,/tikz/every even column/.append style={column sep=6pt}},
  table/col sep=comma]
\addplot[omThm,very thick] table[x=s,y=cb]{figdata/derivatives/21-convertibles-and-the-credit-equity-link/profile.csv};
\addplot[omThm,thick,dashed] table[x=s,y=cb_const]{figdata/derivatives/21-convertibles-and-the-credit-equity-link/profile.csv};
\addplot[omDef,very thick] table[x=s,y=floor]{figdata/derivatives/21-convertibles-and-the-credit-equity-link/profile.csv};
\addplot[omDef,thick,dashed] table[x=s,y=floor_const]{figdata/derivatives/21-convertibles-and-the-credit-equity-link/profile.csv};
\addplot[gray,thick,dotted] table[x=s,y=conv]{figdata/derivatives/21-convertibles-and-the-credit-equity-link/profile.csv};
\legend{convertible (action-crédit),convertible (intensité constante),plancher obligataire (action-crédit),plancher obligataire (intensité constante),valeur de conversion}
\end{axis}
\end{tikzpicture}

\omcaption{La convertible du chapitre (\omterm{def:dv:convertibles-and-the-credit-equity-link:price}{prix de conversion} 40, action à 30 aujourd'hui) selon le cours de l'action, sous une intensité constante et sous une
intensité de défaut action-crédit qui croît quand l'action baisse. Le \omterm{def:dv:convertibles-and-the-credit-equity-link:floor}{plancher obligataire} est plat dans le premier modèle et baisse avec l'action dans le
second~; la convertible suit la \omterm{def:dv:convertibles-and-the-credit-equity-link:value}{valeur de conversion} à droite et le plancher à gauche. Données~: le
tutoriel.}\label{fig:dv:convertibles-and-the-credit-equity-link:profile}
\end{omfigure}

\section{Modèles action-crédit}

La représentation à intensité constante présente un défaut qui compte le plus là où le risque du desk est le plus élevé. Lorsque l'action d'une société chute fortement, son
crédit se dégrade en général aussi, et le \omterm{def:dv:convertibles-and-the-credit-equity-link:floor}{plancher obligataire}, sur lequel repose toute la partie gauche du profil, baisse avec lui. Les modèles qui lient
le défaut au cours de l'action rendent compte de cet effet.

\begin{definition}[Modèle action-crédit]\label{def:dv:convertibles-and-the-credit-equity-link:e2c}
Un \emph{modèle action-crédit}\index{modèle action-crédit} (equity-to-credit model) valorise ensemble les instruments actions et crédit d'un même émetteur en faisant de
l'intensité de défaut une fonction décroissante du cours de l'action, par exemple $\lambda(S)=\lambda_0(S/S_0)^{-p}$. En cas de défaut, le cours
saute à zéro et les obligataires recouvrent une fraction de la valeur nominale.
\end{definition}

Sous la mesure de valorisation, l'action doit tout de même rapporter en moyenne le taux sans risque, si bien que sa dérive est relevée de $\lambda(S)$ pour compenser
le saut à zéro. La valeur $V(t,S)$ de la convertible vérifie, en $x=\ln S$,
\[
\partial_tV+\tfrac12\sigma^2\partial_{xx}V+\bigl(r-q+\lambda(S)-\tfrac12\sigma^2\bigr)\partial_xV-\bigl(r+\lambda(S)\bigr)V+\lambda(S)\,RF=0,
\]
entre les dates de coupon, sous les contraintes suivantes~: $V\ge\kappa S$ pour la conversion, $V\le\max(\text{prix de rachat},\kappa S)$ lorsque le call est exerçable, et
$V\ge\text{prix du put}$ aux dates de put. C'est l'\omterm{def:dv:black-scholes-three-ways:pde}{équation de Black--Scholes} (chapitre~3) avec un taux d'extinction et un terme source de recouvrement. Ayache, Forsyth
et Vetzal exposent cette formulation à une seule équation. Tsiveriotis et Fernandes avaient auparavant scindé la convertible en une composante purement monétaire
actualisée au taux risqué et une composante action actualisée au taux sans risque. La construction résout l'équation par Crank--Nicolson sur une
grille de 500 points en log-prix (différences finies, livre~4, chapitre~27, et chapitre~22 de ce livre). Elle applique les contraintes après
chaque pas, la même projection que celle qui valorisait les options américaines au chapitre~6. Andersen et Buffum ont montré qu'une calibration naïve de tels modèles
peut biaiser sensiblement les prix, et que la fonction d'intensité doit être calibrée conjointement sur la courbe de crédit de l'émetteur et sur ses options.

\begin{omfigure}
\begin{tikzpicture}
\begin{axis}[width=10.5cm,height=4.6cm,
  xlabel={cours de l'action},ylabel={intensité de défaut (\% par an)},
  tick label style={font=\small},label style={font=\small},
  xmin=10,xmax=90,ymin=0,ymax=20,grid=major,grid style={gray!25},
  legend columns=2,legend cell align=left,legend style={font=\footnotesize,at={(0.5,-0.32)},anchor=north,draw=none,fill=none,/tikz/every even column/.append style={column sep=8pt}},
  table/col sep=comma]
\addplot[omThm,very thick] table[x=s,y=hazard]{figdata/derivatives/21-convertibles-and-the-credit-equity-link/profile.csv};
\draw[omDef,thick,dashed] (axis cs:10,5) -- (axis cs:90,5);
\legend{{$\lambda(S)=5\%\,(S/30)^{-1.2}$},constante 5~\%}
\end{axis}
\end{tikzpicture}

\omcaption{Les deux intensités du chapitre, égales au cours actuel de l'action, soit 30 (un spread de crédit de 300 points de base avec un recouvrement de 40~\%).
L'intensité action-crédit vaut 18.7~\% par an à un cours de 10 et 1.5~\% à 80~; la ligne en pointillés est l'intensité constante. Données~:
le code du chapitre.}\label{fig:dv:convertibles-and-the-credit-equity-link:hazard}
\end{omfigure}

\begin{example}[Ce que change le lien avec le crédit]\label{ex:dv:convertibles-and-the-credit-equity-link:e2c}
Prenons $\lambda(S)=5\%\,(S/30)^{-1.2}$, le même spread de crédit de 300 points de base au cours actuel. L'intensité vaut 18.7~\% par an à un cours de 10
(\cref{fig:dv:convertibles-and-the-credit-equity-link:hazard}). Le \omterm{def:dv:convertibles-and-the-credit-equity-link:floor}{plancher
obligataire} vaut désormais 78.1 à 20 et 90.4 à 80, alors qu'il valait 83.2 aux deux cours. L'\omterm{def:dv:convertibles-and-the-credit-equity-link:cb}{obligation convertible} vaut 83.9 à 20 contre 89.1 avec l'intensité constante, et
98.5 à 30 contre 100.2. Le delta vaut 1.55 action par obligation à 30 contre 1.37, et 1.39 contre 0.84 à 20. Une baisse de l'action pèse désormais aussi sur le \omterm{def:dv:convertibles-and-the-credit-equity-link:floor}{plancher
obligataire}, de sorte que la valeur de l'obligation baisse plus vite. Sous un cours d'environ 19, le gamma devient négatif~: $-0.03$ à 15 et $-0.10$ à 10. La convexité que le desk croyait détenir disparaît
là où le crédit prend le relais.
\end{example}

\begin{omfigure}
\begin{tikzpicture}
\begin{groupplot}[group style={group size=2 by 1,horizontal sep=1.6cm},
  width=6.6cm,height=5cm,
  tick label style={font=\small},label style={font=\small},grid=major,grid style={gray!25},table/col sep=comma]
\nextgroupplot[xlabel={cours de l'action},ylabel={delta (actions par obligation)},xmin=10,xmax=90,ymin=0,ymax=2.6,
  title={\small delta},
  legend columns=1,legend cell align=left,legend style={font=\footnotesize,at={(0.5,-0.32)},anchor=north,draw=none,fill=none}]
\addplot[omThm,very thick] table[x=s,y=delta]{figdata/derivatives/21-convertibles-and-the-credit-equity-link/profile.csv};
\addplot[omDef,thick,dashed] table[x=s,y=delta_const]{figdata/derivatives/21-convertibles-and-the-credit-equity-link/profile.csv};
\legend{modèle action-crédit,intensité constante}
\nextgroupplot[xlabel={cours de l'action},ylabel={gamma},xmin=10,xmax=90,ymin=-0.12,ymax=0.08,
  yticklabel style={/pgf/number format/fixed,/pgf/number format/precision=2},
  title={\small gamma}]
\addplot[omThm,very thick] table[x=s,y=gamma]{figdata/derivatives/21-convertibles-and-the-credit-equity-link/profile.csv};
\addplot[omDef,thick,dashed] table[x=s,y=gamma_const]{figdata/derivatives/21-convertibles-and-the-credit-equity-link/profile.csv};
\draw[gray] (axis cs:10,0) -- (axis cs:90,0);
\end{groupplot}
\end{tikzpicture}

\omcaption{Le delta et le gamma de l'\omterm{def:dv:convertibles-and-the-credit-equity-link:cb}{obligation convertible} selon les deux modèles. Avec le lien de crédit, le delta reste élevé quand l'action baisse, car le
\omterm{def:dv:convertibles-and-the-credit-equity-link:floor}{plancher obligataire} baisse avec elle, et le gamma devient négatif sous environ 19~: la position à convexité longue sur laquelle repose l'arbitrage devient une position à convexité courte
en situation de détresse. Données~: le tutoriel.}\label{fig:dv:convertibles-and-the-credit-equity-link:greeks}
\end{omfigure}

\section{Ce que couvre un desk de convertibles}

La position classique est longue de l'\omterm{def:dv:convertibles-and-the-credit-equity-link:cb}{obligation convertible} et vendeuse à découvert du delta en actions. Couverte ainsi, elle profite d'une volatilité réalisée supérieure à la volatilité
implicite contenue dans son prix (chapitre~4), comme le ferait une option achetée. Elle conserve le risque de crédit du \omterm{def:dv:convertibles-and-the-credit-equity-link:floor}{plancher obligataire} et le risque du delta du modèle.
Les desks couvrent le crédit au moyen d'un credit default swap (CDS) (livre~2, chapitre~23), dimensionné sur la sensibilité au crédit de l'obligation, et ils redoutent surtout
le scénario combiné.

\begin{definition}[Arbitrage de convertibles]\label{def:dv:convertibles-and-the-credit-equity-link:arb}
L'\emph{arbitrage de convertibles}\index{arbitrage de convertibles} est la stratégie qui consiste à acheter des \omterm{def:dv:convertibles-and-the-credit-equity-link:cb}{obligations convertibles} cotées sous leur valeur de modèle,
à vendre à découvert les actions de l'émetteur à hauteur du delta du modèle et à acheter souvent une protection de crédit, afin de gagner la convexité de l'obligation, son coupon
net du coût de la position courte, et la convergence de son prix vers sa valeur.
\end{definition}

\begin{example}[Le portefeuille de couverture]\label{ex:dv:convertibles-and-the-credit-equity-link:hedge}
Au cours de 30, le delta du \omterm{def:dv:convertibles-and-the-credit-equity-link:e2c}{modèle action-crédit} vaut 1.55 action par obligation. Un élargissement de 1 point de base du spread de crédit de l'émetteur, le cours de l'action restant fixe,
coûte à l'\omterm{def:dv:convertibles-and-the-credit-equity-link:cb}{obligation convertible} 0.0167 par obligation (1.67 pour 100 points de base), soit environ la moitié des 0.0296 du \omterm{def:dv:convertibles-and-the-credit-equity-link:floor}{plancher obligataire}, car le droit de conversion
en amortit l'effet. Avec une annuité risquée à cinq ans de 4.12, une protection de crédit sur 40.6~\% du nominal compense cette sensibilité.
\end{example}

\begin{example}[L'action baisse de 20~\%, spread de crédit de 300 points de base plus large]\label{ex:dv:convertibles-and-the-credit-equity-link:stress}
L'action passe de 30 à 24, et le spread de crédit de l'émetteur à ce cours finit 300 points de base au-dessus de celui d'aujourd'hui, soit une intensité de 10~\%. L'\omterm{def:dv:convertibles-and-the-credit-equity-link:cb}{obligation convertible} passe de
98.47 à 86.62, soit une perte de 11.84 pour 100 de nominal. La position vendeuse de 1.55 action gagne 9.30. La position couverte en delta perd 2.55 pour 100 de nominal, soit 2.55
millions sur 100 millions. Couverte avec le delta du modèle à intensité constante, de 1.37, elle perd 4.38. La protection de crédit gagne 4.47 dans le même
scénario et transforme la perte en gain de 1.93. Le seul élargissement du spread de crédit, sans mouvement de l'action, coûte à la position couverte en delta 3.35. Le mouvement de l'action
seul, le spread de crédit restant inchangé, lui rapporte 2.21 de convexité (\cref{fig:dv:convertibles-and-the-credit-equity-link:stress}).
\end{example}

\begin{omfigure}
\begin{tikzpicture}
\begin{axis}[width=10.5cm,height=5cm,
  xlabel={mouvement de l'action (\%)},ylabel={P\&L couvert en delta (pour 100 de nominal)},
  tick label style={font=\small},label style={font=\small},
  xmin=-32,xmax=22,ymin=-4,ymax=4.5,grid=major,grid style={gray!25},
  legend columns=3,legend cell align=left,legend style={font=\footnotesize,at={(0.5,-0.3)},anchor=north,draw=none,fill=none,/tikz/every even column/.append style={column sep=8pt}},
  table/col sep=comma]
\addplot[omThm,very thick,mark=*,mark size=1.4pt] table[x=move,y=w0]{figdata/derivatives/21-convertibles-and-the-credit-equity-link/stress.csv};
\addplot[omDef,very thick,mark=square*,mark size=1.4pt] table[x=move,y=w150]{figdata/derivatives/21-convertibles-and-the-credit-equity-link/stress.csv};
\addplot[omProp,very thick,mark=triangle*,mark size=1.6pt] table[x=move,y=w300]{figdata/derivatives/21-convertibles-and-the-credit-equity-link/stress.csv};
\draw[gray] (axis cs:-32,0) -- (axis cs:22,0);
\legend{spread inchangé,+150 pb,+300 pb}
\end{axis}
\end{tikzpicture}

\omcaption{\omterm{def:dv:convertibles-and-the-credit-equity-link:cb}{Obligation convertible} longue, vendeuse de 1.55 action par obligation (delta du \omterm{def:dv:convertibles-and-the-credit-equity-link:e2c}{modèle action-crédit} au cours de 30), après un mouvement de l'action et un
élargissement du spread de crédit de l'émetteur au nouveau cours de l'action. La convexité rapporte lorsque le spread tient~; un élargissement transforme chaque scénario en perte. Données~:
le tutoriel.}\label{fig:dv:convertibles-and-the-credit-equity-link:stress}
\end{omfigure}

Le choix du delta par le modèle est lui-même une décision de couverture. Un delta action-crédit est plus grand, car il tient compte du canal de crédit par lequel une
baisse de l'action abaisse le \omterm{def:dv:convertibles-and-the-credit-equity-link:floor}{plancher obligataire}, et il protège donc davantage dans le scénario défavorable. Un desk qui utilise le delta à intensité constante est implicitement
long en crédit à travers sa couverture en actions. Il est aussi exposé au saut de défaut, que le delta et un CDS dimensionnés pour de petits mouvements ne couvrent qu'en partie
(livre~6, chapitre~13, et les modèles structurels de son chapitre~14).

\section{Arbitrage de convertibles et ses crises}

L'arbitrage tire son rendement de l'écart entre la valeur donnée par un modèle et un cours de marché, et de la convexité qu'il détient. Il est financé~:
le fonds emprunte en nantissant les convertibles et prête les actions qu'il vend par l'intermédiaire de son prime broker. Ce financement constitue le point faible de l'arbitrage. Lorsque
de nombreux fonds détiennent les mêmes obligations et doivent vendre au même moment, les cours baissent ensemble, à l'écart de toute valeur de modèle. Les couvertures protègent alors
contre le mauvais risque.

Le «~Slow moving capital~» de Mitchell, Pedersen et Pulvino (2007), ainsi que l'étude de Mitchell et Pulvino sur 2008 (2012), examinent ce qui se produit alors.
Selon leur analyse, les hedge funds détenaient jusqu'à 75~\% du marché des convertibles~; au début de 2005, de gros investisseurs ont commencé à retirer leurs capitaux
des fonds d'\omterm{def:dv:convertibles-and-the-credit-equity-link:arb}{arbitrage de convertibles}, dont plus de 20~\% au premier trimestre (un chiffre qu'ils tirent du Barclay Group), et les ventes forcées
des fonds ont fait passer les convertibles sous leurs valeurs fondamentales. Pour 2008, Mitchell et Pulvino décrivent le mécanisme. La défaillance imminente de grands prime brokers a brusquement réduit l'effet de levier disponible pour les hedge funds,
en apparence, un financement de long terme est devenu de court terme, et les fonds de valeur relative n'ont plus pu détenir des actifs semblables à des prix semblables. Pour un modèle, cela
désigne un scénario que l'exemple de stress ci-dessus ne contient pas~: la convertible se déprécie par rapport à sa propre valeur de modèle, tandis que l'action et le crédit
bougent à peine. Aucune couverture sur l'action ou sur le CDS ne la couvre. Sa protection réside dans un levier plus faible et dans un financement qui ne peut pas être retiré du jour au lendemain.

\section{Tutoriel~: une convertible sur une grille}\label{tut:dv:convertibles-and-the-credit-equity-link:tut}

\begin{tutorial}
\textbf{Objectif.} Évaluer la convertible du chapitre sous un taux de défaut constant puis sous un taux de défaut action-crédit, calculer son delta, son gamma et sa sensibilité au crédit,
et soumettre la position couverte à un stress. \textbf{État final~:} les quatre chiffres et les nombres du problème du week-end.
\begin{tutsteps}
\item \textbf{L'opérateur}~: dérive relevée par le taux de défaut, extinction à $r+\lambda$, recouvrement comme source~:
\omcode{code/firm/convertible/firm_convertible.py}{70}{74}{L'opérateur de valorisation de la convertible.}\label{lst:dv:convertibles-and-the-credit-equity-link:operator}
\item \textbf{Les contraintes après chaque pas}~: coupons, options de vente, conversion et \omterm{def:dv:convertibles-and-the-credit-equity-link:softcall}{soft call}~:
\omcode{code/firm/convertible/firm_convertible.py}{92}{100}{Coupons, options de vente, conversion et rachat.}\label{lst:dv:convertibles-and-the-credit-equity-link:constraints}
\item \textbf{Exécuter} \texttt{dv\_convertible.table()}, \texttt{credit\_sensitivity()}, \texttt{stress(-0.20, 300)}, \texttt{call\_effect()} et
\texttt{fig\_convertible.py}.
\end{tutsteps}
\textbf{Pour aller plus loin.} Poser $p=2$ et recalculer le delta à 20~; ajouter une option de vente investisseur à 100 la troisième année et observer le déplacement du \omterm{def:dv:convertibles-and-the-credit-equity-link:floor}{plancher obligataire}~; valoriser
l'option de rachat dans l'hypothèse d'un rachat ferme (inconditionnel) à partir de la deuxième année.
\end{tutorial}

\section{Construction~: le pricer de convertibles}\label{bld:dv:convertibles-and-the-credit-equity-link:convertible}

\begin{build}
\bfield{Rôle} Le moteur de valorisation et de risque des convertibles de la firme miniature~: valeur, delta, gamma et sensibilité au crédit de chaque obligation de son
carnet de convertibles, sous un taux de défaut pouvant dépendre de l'action.
\bfield{Interface} \texttt{Convertible(face, maturity, coupon, freq, ratio, recovery, call\_start, call\_price, call\_trigger, puts)}~;
\texttt{power\_hazard(lam0, s0, p, cap)}~; \texttt{price\_grid(cb, r, q, vol, hazard, s\_min, s\_max, nx, steps\_per\_year) -> (spots, values)}~;
\texttt{value\_at}~; \texttt{greeks(s0, grid) -> value, delta, gamma}~; \texttt{bond\_floor(cb, r, spread)}.
\bfield{Règles} Le taux de défaut est calibré sur la courbe de crédit de l'émetteur avant que toute convertible ne soit valorisée~; la conversion, le rachat et l'option de vente sont appliqués
après chaque pas de temps~; la sensibilité au crédit est présentée avec le delta action, et le stress de l'action et du spread conjointement est exécuté chaque jour.
\bfield{Tests d'acceptation} \texttt{code/firm/convertible/tests/}~: sans conversion ni taux de défaut, le prix est celui de l'obligation~; avec un taux de défaut constant, la
formule fermée avec recouvrement du nominal~; la valeur dépasse le \omterm{def:dv:convertibles-and-the-credit-equity-link:floor}{plancher obligataire} et la \omterm{def:dv:convertibles-and-the-credit-equity-link:value}{valeur de conversion} et croît avec l'action~; le rachat abaisse et l'option de vente relève la valeur~; le
lien de crédit relève le delta aux faibles cours de l'action.
\bfield{Pour aller plus loin} La décomposition de Tsiveriotis--Fernandes~; un taux de défaut calibré sur une courbe de CDS et sur les options de l'émetteur~; les dividendes et les préavis
de rachat~; les convertibles conditionnelles du livre~2, chapitre~26.
\end{build}

\begin{omsources}
\item K.~Tsiveriotis et C.~Fernandes, ``Valuing convertible bonds with credit risk'', \emph{Journal of Fixed Income} 8(2) (1998) 95--102.
\item E.~Ayache, P.~A.~Forsyth et K.~R.~Vetzal, ``Valuation of convertible bonds with credit risk'', \emph{Journal of Derivatives} 11(1) (2003)
9--29.
\item L.~Andersen et D.~Buffum, ``Calibration and implementation of convertible bond models'', \emph{Journal of Computational Finance} 7 (2003)
1--34.
\item M.~Mitchell, L.~H.~Pedersen et T.~Pulvino, ``Slow moving capital'', \emph{American Economic Review} 97(2) (2007) 215--220.
\item M.~Mitchell et T.~Pulvino, ``Arbitrage crashes and the speed of capital'', \emph{Journal of Financial Economics} 104(3) (2012) 469--490.
\end{omsources}

\section{Exercices}

\begin{exercise}[$\star$]\label{exo:dv:convertibles-and-the-credit-equity-link:1}
Une convertible de nominal 100 et de \omterm{def:dv:convertibles-and-the-credit-equity-link:ratio}{ratio de conversion} 2.5 se traite à 110 alors que l'action vaut 36. Donner le \omterm{def:dv:convertibles-and-the-credit-equity-link:price}{prix de conversion}, la \omterm{def:dv:convertibles-and-the-credit-equity-link:value}{valeur de conversion} et la
\omterm{def:dv:convertibles-and-the-credit-equity-link:premium}{prime de conversion}.
\end{exercise}

\begin{exercise}[$\star$]\label{exo:dv:convertibles-and-the-credit-equity-link:2}
Pourquoi une convertible vaut-elle au moins son \omterm{def:dv:convertibles-and-the-credit-equity-link:floor}{plancher obligataire} et au moins sa \omterm{def:dv:convertibles-and-the-credit-equity-link:value}{valeur de conversion}~?
\end{exercise}

\begin{exercise}[$\star$]\label{exo:dv:convertibles-and-the-credit-equity-link:3}
Pourquoi le \omterm{def:dv:convertibles-and-the-credit-equity-link:softcall}{soft call} diminue-t-il la valeur de la convertible, et où l'effet est-il le plus fort~?
\end{exercise}

\begin{exercise}[$\star\star$]\label{exo:dv:convertibles-and-the-credit-equity-link:4}
Expliquer pourquoi la dérive de l'action sous la mesure de valorisation vaut $r-q+\lambda(S)$ dans un \omterm{def:dv:convertibles-and-the-credit-equity-link:e2c}{modèle action-crédit}.
\end{exercise}

\begin{exercise}[$\star\star$]\label{exo:dv:convertibles-and-the-credit-equity-link:5}
Pourquoi la sensibilité au crédit de la convertible est-elle plus faible que celle de son \omterm{def:dv:convertibles-and-the-credit-equity-link:floor}{plancher obligataire}~?
\end{exercise}

\begin{exercise}[$\star\star$]\label{exo:dv:convertibles-and-the-credit-equity-link:6}
Expliquer pourquoi le delta action-crédit est plus grand que le delta à taux de défaut constant aux faibles cours de l'action, et ce que cela implique pour une couverture.
\end{exercise}

\begin{exercise}[$\star\star\star$]\label{exo:dv:convertibles-and-the-credit-equity-link:7}
\emph{Programmation.} Recalculer le stress de l'\cref{ex:dv:convertibles-and-the-credit-equity-link:stress} avec l'exposant action-crédit $p=2$ (même
spread aujourd'hui). La perte de la position couverte en delta augmente-t-elle ou diminue-t-elle~?
\end{exercise}

\begin{exercise}[$\star\star\star$]\label{exo:dv:convertibles-and-the-credit-equity-link:8}
\emph{Repérer l'erreur.} «~Notre carnet de convertibles est couvert en delta et par CDS, donc son seul risque est la volatilité.~»
\end{exercise}

\section{Problème~: action en baisse, crédit en hausse}

\begin{problem}[{Problème du week-end --- le scénario que redoute un desk de convertibles}]\label{pb:dv:convertibles-and-the-credit-equity-link:1}
Un desk détient un nominal de 100 millions de la convertible du chapitre (action à 30, \omterm{def:dv:convertibles-and-the-credit-equity-link:price}{prix de conversion} 40, spread de 300 points de base), et il est vendeur
en actions, du delta action-crédit.

\medskip\noindent\textbf{Partie I --- L'obligation.}
\begin{enumerate}
\item Donner la valeur de la convertible, son \omterm{def:dv:convertibles-and-the-credit-equity-link:floor}{plancher obligataire}, sa \omterm{def:dv:convertibles-and-the-credit-equity-link:value}{valeur de conversion} et sa \omterm{def:dv:convertibles-and-the-credit-equity-link:premium}{prime de conversion} à 30 sous le \omterm{def:dv:convertibles-and-the-credit-equity-link:e2c}{modèle action-crédit}.
\item Donner le delta et le nombre d'actions que le desk est vendeur.
\item Donner la valeur à un cours de l'action de 10 et le taux de défaut en ce point.
\item Que coûte le \omterm{def:dv:convertibles-and-the-credit-equity-link:softcall}{soft call} au porteur à un cours de l'action de 52~?
\item Pourquoi le \omterm{def:dv:convertibles-and-the-credit-equity-link:floor}{plancher obligataire} n'est-il pas plat dans ce modèle~?
\end{enumerate}

\medskip\noindent\textbf{Partie II --- Les couvertures.}
\begin{enumerate}[resume]
\item Donner la sensibilité au crédit par point de base, pour la position.
\item Dimensionner la protection en CDS qui la compense.
\item Que gagne la position si l'action baisse de 20~\% et que le spread reste inchangé~?
\item Que perd-elle si le spread s'élargit de 300 points de base et que l'action reste inchangée~?
\item Quel serait le delta du modèle à taux de défaut constant, et de combien d'actions en moins le desk serait-il vendeur~?
\end{enumerate}

\medskip\noindent\textbf{Partie III --- Le scénario.}
\begin{enumerate}[resume]
\item L'action baisse de 20~\% et le spread de crédit, au nouveau prix, s'élargit finalement de 300 points de base. Donner la nouvelle valeur de l'\omterm{def:dv:convertibles-and-the-credit-equity-link:cb}{obligation convertible} et la perte de la position couverte en delta.
\item Donner la perte avec le delta du modèle à taux de défaut constant.
\item Donner le résultat en ajoutant la couverture en CDS.
\item Déterminer la part de la perte qui relève du crédit et celle qui relève du delta du modèle.
\item Identifier le scénario qui manque à cette analyse.
\end{enumerate}

\medskip\noindent\textbf{Partie IV --- Jugement.}
\begin{enumerate}[resume]
\item Choisir le delta sur lequel opérer, et justifier ce choix.
\item Déterminer la quantité de CDS à détenir, compte tenu du saut au défaut.
\item Préciser les limites à fixer pour le carnet.
\item Énoncer le \emph{résultat de référence}~: la perte de la position d'arbitrage sur \omterm{def:dv:convertibles-and-the-credit-equity-link:cb}{obligations convertibles} couverte en delta lorsque le spread de crédit de l'émetteur s'élargit de 300 points
de base pendant que son action baisse de 20~\%.
\item En une phrase, dire ce que détient un arbitragiste sur \omterm{def:dv:convertibles-and-the-credit-equity-link:cb}{obligations convertibles}.
\end{enumerate}
\end{problem}

\section{Questions d'entretien}

\begin{interviewq}[$\star$ \iqroles{trader, chercheur}]\label{iq:dv:convertibles-and-the-credit-equity-link:1}
Esquisser la valeur d'une \omterm{def:dv:convertibles-and-the-credit-equity-link:cb}{obligation convertible} en fonction du cours de l'action. Nommer les régions.
\end{interviewq}

\begin{interviewq}[$\star\star$ \iqroles{chercheur}]\label{iq:dv:convertibles-and-the-credit-equity-link:2}
Écrire l'équation de prix d'une \omterm{def:dv:convertibles-and-the-credit-equity-link:cb}{obligation convertible} soumise au risque de défaut. Quelles sont les conditions aux limites et les contraintes~?
\end{interviewq}

\begin{interviewq}[$\star\star$ \iqroles{trader}]\label{iq:dv:convertibles-and-the-credit-equity-link:3}
Que gagne un arbitragiste sur \omterm{def:dv:convertibles-and-the-credit-equity-link:cb}{obligations convertibles}, et quels sont les risques~?
\end{interviewq}

\begin{interviewq}[$\star\star$ \iqroles{développeur}]\label{iq:dv:convertibles-and-the-credit-equity-link:4}
Comment implémenteriez-vous les clauses de conversion, de rachat et de vente dans un pricer par différences finies~?
\end{interviewq}

\begin{interviewq}[$\star\star$ \iqroles{risque}]\label{iq:dv:convertibles-and-the-credit-equity-link:5}
Concevoir un test de résistance pour un carnet d'arbitrage sur \omterm{def:dv:convertibles-and-the-credit-equity-link:cb}{obligations convertibles}.
\end{interviewq}

\begin{interviewq}[$\star\star\star$ \iqroles{trader, risque}]\label{iq:dv:convertibles-and-the-credit-equity-link:6}
Les prime brokers réduisent de moitié l'effet de levier du carnet du jour au lendemain. Que se passe-t-il pour les positions et pour leurs prix~?
\end{interviewq}
